%%%%%%%%%%  scan idx 51  |  folio 40 (continued)  %%%%%%%%%%

Sea $P:\underline A\longrightarrow\underline C_\Lambda$ una
categoría cofibrada en grupoides. Esta da entonces una categoría
cofibrada en grupoides
\[
 \operatorname{pro}(P):\operatorname{pro}(\underline A)
 \longrightarrow\operatorname{pro}(\underline C_\Lambda)
\]
(para la definición de la categoría de proobjetos
$\operatorname{pro}(\underline A)$, véase A. Grothendieck,
S\'eminaire Bourbaki 195), que encaja en el diagrama conmutativo
\[
\begin{tikzcd}[column sep=large,row sep=large]
\underline A \arrow[r] \arrow[d,"P"']
 & \operatorname{pro}(\underline A)
   \arrow[d,"\operatorname{pro}(P)"] \\
\underline C_\Lambda \arrow[r]
 & \operatorname{pro}(\underline C_\Lambda).
\end{tikzcd}
\]

Sea $\widehat{\underline C}_\Lambda$ la subcategoría plena de
$\operatorname{pro}(\underline C_\Lambda)$ cuyos objetos son sistemas
$(R_i)_{i\in\mathbb N}$ con $R_i\in\underline C_\Lambda$ tales que
cada morfismo $\varphi_i:R_i\longrightarrow R_{i-1}$ es sobreyectivo
y, para todo $i$ suficientemente grande, induce un isomorfismo
\[
 \frac{m_i}{m_\Lambda R_i+m_i^2}
 \xrightarrow{\ \sim\ }
 \frac{m_{i-1}}{m_\Lambda R_{i-1}+m_{i-1}^2},
\]
donde $\mathbb N$ denota la categoría ordenada de los números naturales
y $m_i$ es el ideal maximal de $R_i$.

Sea $\widehat{\underline A}$ la subcategoría plena de
$\operatorname{pro}(\underline A)$ definida por
\[
 \operatorname{ob}(\widehat{\underline A})=
 \left\{a\in\operatorname{ob}(\operatorname{pro}(\underline A))
 \mathrel{\Big|}
 \operatorname{pro}(P)(a)\in
 \operatorname{ob}(\widehat{\underline C}_\Lambda)\right\}.
\]
Obtenemos una categoría cofibrada en grupoides
$\widehat P:\widehat{\underline A}\longrightarrow
\widehat{\underline C}_\Lambda$ y un diagrama conmutativo
\[
\begin{tikzcd}[column sep=large,row sep=large]
\underline A \arrow[r,hook] \arrow[d,"P"']
 & \widehat{\underline A} \arrow[d,"\widehat P"] \\
\underline C_\Lambda \arrow[r,hook]
 & \widehat{\underline C}_\Lambda.
\end{tikzcd}
\]
La categoría $\widehat{\underline C}_\Lambda$ es canónicamente
equivalente a la categoría de las $\Lambda$-álgebras locales completas
de tipo formalmente finito sobre $\Lambda$ con cuerpo residual fijo
$k$; en adelante efectuaremos esta identificación.

\underline{Definición 1.9. Sea $\underline A$ un grupoide cofibrado
sobre $\underline C_\Lambda$. Un objeto $\xi$ de
$\widehat{\underline A}$ se llama cuasiinicial si, para todo objeto
$a$ de $\widehat{\underline A}$, existe un morfismo $\xi\to a$ en
$\widehat{\underline A}$. Un objeto cuasiinicial $\xi$ se llama
mínimo si el morfismo
$\xi:h_R\longrightarrow[\widehat{\underline A}]$ induce una biyección}

%%%%%%%%%%  scan idx 52  |  folio 41  %%%%%%%%%%

\[
 \xi(k[\varepsilon]):
 \operatorname{Hom}_{\Lambda\text{-alg}}(R,k[\varepsilon])
 \xrightarrow{\ \sim\ }
 [\underline A(k[\varepsilon])],
 \qquad R=P(\xi).
\]
\underline{Un objeto $\xi$ de $\widehat{\underline A}$ se llama
objeto formal versal de $\underline A$ si todo morfismo sobreyectivo
$\alpha:a'\to a$ en $\underline A$ induce un morfismo sobreyectivo}
\[
 \operatorname{Hom}_{\widehat{\underline A}}(\xi,a')
 \longrightarrow
 \operatorname{Hom}_{\widehat{\underline A}}(\xi,a).
\]
\underline{Un objeto formal versal $\xi$ se llama mínimo si
$\xi:h_R\longrightarrow[\widehat{\underline A}]$ induce la biyección
anterior sobre $k[\varepsilon]$, donde $R=P(\xi)$.}

\underline{Proposición 1.10. Sea $\underline A$ un grupoide cofibrado
sobre $\underline C_\Lambda$.}

(i) \underline{Todo objeto formal (mínimamente) versal de
$\underline A$ es un objeto cuasiinicial (mínimo) de
$\widehat{\underline A}$.}

(ii) \underline{Supongamos que $\xi$ es un objeto de
$\sisi{\widehat A}{\widehat{\underline A}}$ para el cual}
\[
 \widetilde\xi(k[\varepsilon]):
 \operatorname{Hom}_{\Lambda\text{-alg}}(R,k[\varepsilon])
 \longrightarrow[\underline A(k[\varepsilon])]
\]
\underline{es inyectivo, donde $P(\xi)=R$. Entonces todo endomorfismo
de $\xi$ en $\sisi{\widehat A}{\widehat{\underline A}}$ es un automorfismo. En
particular, un objeto cuasiinicial mínimo de $\widehat{\underline A}$,
si existe, es único salvo isomorfismo no canónico.}

\underline{Demostración.} (i) Sea $\xi$ un objeto formal versal de
$\underline A$. Para un objeto $a$ de $\widehat{\underline A}$,
elijamos una presentación mediante morfismos sobreyectivos en
$\underline A$,
\[
 \cdots\longrightarrow a_n\xrightarrow{\alpha_n}a_{n-1}
 \longrightarrow\cdots\longrightarrow
 a_1\xrightarrow{\alpha_1}a_0,
 \qquad a=\varprojlim_n a_n.
\]
Supongamos elegidos morfismos $\varphi_i:\xi\longrightarrow a_i$ con
$\alpha_i\varphi_i=\varphi_{i-1}$ para todo $i<n$. Como $\xi$ es
versal y $\alpha_n$ es sobreyectivo, existe un morfismo
$\varphi_n:\xi\longrightarrow a_n$ con
$\alpha_n\varphi_n=\varphi_{n-1}$. Así,
$\varphi=\varprojlim_n\varphi_n$ define un morfismo
$\xi\longrightarrow a$.

(ii) Sea $\alpha:\xi\longrightarrow\xi$ un endomorfismo en
$\widehat{\underline A}$. Para todo
$\lambda\in\operatorname{Hom}_{\Lambda\text{-alg}}
(R,k[\varepsilon])$, se tiene
\[
 \lambda_*\xi\simeq
 (\lambda\circ P(\alpha))_*\xi.
\]
La inyectividad da, por tanto,
$\lambda=\lambda\circ P(\alpha)$ para todo $\lambda$. De ello se
deduce que $P(\alpha):R\longrightarrow R$ es sobreyectivo. Todo
endomorfismo sobreyectivo de un anillo noetheriano es un automorfismo;
por consiguiente, $P(\alpha)$ y, por ende, $\alpha$ son automorfismos.

